\documentclass[11pt]{article}
\usepackage[margin=1in]{geometry}
\usepackage{amsmath,amssymb,mathtools}
\usepackage[hidelinks]{hyperref}
\usepackage{microtype}

\newcommand{\Cond}{\operatorname{Cond}}
\newcommand{\lmax}{\lambda_{\max}}
\newcommand{\lmin}{\lambda_{\min}}

\title{NERF Erasure-Rate Monotonicity\\
\large Literature-Already-Answered Identification}
\author{Run \texttt{fa\_banach\_001}, lane 2}
\date{9 August 2026}

\begin{document}
\maketitle

\noindent\textbf{Status.}
\texttt{literature\_already\_answered}. This note records an exact later
answer; it is not claimed as a new discovery.

\section{Original question}

Let $F=[f_1\ \cdots\ f_N]$ be an $M\times N$ frame. Fickus and
Mixon call $F$ a $(p,C)$-numerically erasure-robust frame (NERF) if every
submatrix $F_{\mathcal K}$ with
$|\mathcal K|=K=(1-p)N$ satisfies
$\Cond(F_{\mathcal K})\leq C$.

In Section 4, source-PDF page 14, they ask whether a $(p,C)$-NERF must
also be a $(p',C)$-NERF for every $p'\in[0,p)$ for which
$L=(1-p')N$ is integral \cite{fickus-mixon}. Equivalently, does the same
worst-case condition-number bound persist when more columns survive?

\section{Exact later answer}

Liu explicitly cites the Fickus--Mixon question in the introduction of
\cite{liu}. Lemma 1 (pp.~586--587) proves
\[
 \max_{|\mathcal K|=K}\Cond(F_{\mathcal K})
 \;\geq\;
 \max_{|\mathcal S|=L}\Cond(F_{\mathcal S})
 \qquad (K<L).
\]
Therefore every $(p,C)$-NERF is a $(p',C)$-NERF whenever $p'<p$ and the
two prescribed cardinalities are integral. This is exactly the affirmative
answer requested in the source.

\section{Independent verification}

For completeness, fix a set $\mathcal S$ of $L$ columns and write
$A_{\mathcal T}=F_{\mathcal T}F_{\mathcal T}^{*}$ for each
$K$-element subset $\mathcal T\subset\mathcal S$. If
$\Cond(F_{\mathcal T})\leq C$, then
\[
 \lmax(A_{\mathcal T})\leq C^2\lmin(A_{\mathcal T}).
\]
The cone defined by this inequality is closed under finite sums: for any such
positive semidefinite matrices $A_j$,
\begin{align*}
 \lmax\!\left(\sum_j A_j\right)
 &\leq \sum_j\lmax(A_j)
 \leq C^2\sum_j\lmin(A_j)\\
 &\leq C^2\lmin\!\left(\sum_j A_j\right).
\end{align*}
Every column of $F_{\mathcal S}$ occurs in exactly
$\binom{L-1}{K-1}$ of the $K$-subsets, and hence
\[
 \sum_{\substack{\mathcal T\subset\mathcal S\\|\mathcal T|=K}}
 A_{\mathcal T}
 =\binom{L-1}{K-1}F_{\mathcal S}F_{\mathcal S}^{*}.
\]
The preceding cone inequality and invariance under positive scalar multiples
give
$\Cond(F_{\mathcal S})^2\leq C^2$. Since $\mathcal S$ was arbitrary,
$F$ is a $(p',C)$-NERF. This also independently checks the decisive content
of Liu's Lemma 1.

\section{\hspace{0.35em}Scope and search record}

The identification resolves the monotonicity question only. It does not
address the source's separate questions asking whether NERFs exist for
$p\geq\tfrac12$ or what the largest admissible erasure rate is.

A bounded novelty search used the exact source title, the phrases
``deleting columns does not necessarily worsen'' and
``$(p',C)$-NERF'', and NERF/erasure-rate monotonicity terminology. It found
Liu's 2020 paper, whose introduction explicitly says it answers the
Fickus--Mixon question and whose Lemma 1 gives the exact comparison.

\begin{thebibliography}{9}
\bibitem{fickus-mixon}
M.~Fickus and D.~G. Mixon,
\emph{Numerically erasure-robust frames},
Linear Algebra and its Applications 437 (2012), 1394--1407;
\href{https://arxiv.org/abs/1202.4525}{arXiv:1202.4525}.

\bibitem{liu}
Y.~Liu,
\emph{Comparison on the robustness against erasure rates of numerically
erasure-robust frames},
International Journal of Applied Mathematics 33 (2020), 585--590,
\href{https://doi.org/10.12732/ijam.v33i4.3}{doi:10.12732/ijam.v33i4.3}.
\end{thebibliography}

\end{document}
